% !TEX root = main.tex
\clearpage
\section{Code guide: BachattLM model and trainer}

\texttt{model.py} is a compact executable specification of a modern
decoder-only Transformer. \texttt{train.py} supplies deterministic batches,
optimization, evaluation, logging, and checkpoint resumption. Together they
contain 384 lines of Python, small enough to audit before moving to a larger
distributed framework.

\subsection{\texttt{model.py}}

Now we leave data engineering and enter the model. This defines the actual
decoder-only Transformer.

The architecture is:
\begin{lstlisting}[style=shadedpython,language={},basicstyle=\scriptsize\ttfamily]
token IDs
   |
   v
token embedding
   |
   v
Transformer Block x N
   |
   |-- RMSNorm
   |-- GQA attention
   |    + RoPE
   |-- residual
   |-- RMSNorm
   |-- SwiGLU FFN
   `-- residual
   |
   v
final RMSNorm
   |
   v
language model head
   |
   v
vocabulary logits
\end{lstlisting}

\subsection{Configuration and tensor dimensions}

\texttt{ModelConfig} defines:
\begin{center}
  \renewcommand{\arraystretch}{1.25}
  \small
  \begin{tabular}{@{}l l p{0.55\textwidth}@{}}
    \textbf{Field} & \textbf{Symbol} & \textbf{Meaning} \\
    \hline
    \texttt{vocab\_size} & \(V\) &
    number of distinct tokens the model can emit; size of the embedding
    table and of the final logit vector \\
    \texttt{dim} & \(C\) &
    residual-stream width --- every token is a \(C\)-dimensional vector
    throughout the network \\
    \texttt{n\_layers} & \(L\) &
    how many Transformer blocks are stacked (each block =
    attention + SwiGLU FFN) \\
    \texttt{n\_heads} & \(H\) &
    number of query heads in attention; each head has width
    \(D=C/H\) \\
    \texttt{n\_kv\_heads} & \(H_{kv}\) &
    number of shared key/value heads under GQA (\(H\) must be divisible
    by \(H_{kv}\); each KV head serves \(H/H_{kv}\) query heads) \\
    \texttt{max\_seq\_len} & \(T\) &
    longest context the RoPE cache and causal mask are built for
  \end{tabular}
\end{center}

Example:
\begin{lstlisting}[style=shadedpython]
vocab_size = 4096
dim        = 512
layers     = 8
query heads = 8
KV heads    = 2
sequence    = 1024
\end{lstlisting}
The per-head dimension is
\[
  D=\frac{C}{H}=\frac{\texttt{dim}}{\texttt{n\_heads}}.
\]
So \(512/8=64\): each query head has \(D=64\).

The code asserts that \(C\) is divisible by \(H\) and that \(H\) is divisible
by \(H_{kv}\). It should also assert that \(D\) is even, because RoPE splits
each head into two equal halves.

\subsection{RMSNorm}

Previously we learned LayerNorm. RMSNorm is simpler: it normalizes based on
root mean square,
\[
  \frac{x}{\sqrt{\operatorname{mean}(x^2)+\epsilon}},
\]
then multiplies by a learned weight. Unlike LayerNorm, it does not subtract
the mean first --- it only rescales:

\[
  \operatorname{RMSNorm}(x)
  =g\odot\frac{x}{\sqrt{\operatorname{mean}(x^2)+\epsilon}}.
\]

The implementation computes the normalization in FP32 and casts back to the
activation dtype. This protects the reduction from lower-precision rounding
while preserving mixed-precision throughput in later operations.

\subsection{Rotary position embeddings (RoPE)}

Old-style positional embeddings:
\begin{lstlisting}[style=shadedpython]
token_embedding + position_embedding
\end{lstlisting}
versus RoPE:
\begin{lstlisting}[style=shadedpython,language={},basicstyle=\scriptsize\ttfamily]
token_embedding
      |
      v
Transformer
      |
      v
rotate Q and K based on position
      |
      v
attention
\end{lstlisting}
RoPE is attractive because it gives attention a strong representation of
relative position, works well with long contexts, and has become very common
in modern decoder-only LLMs.

RoPE is a way for a Transformer to understand token position and relative
distance without adding a separate positional vector to each token embedding.

In self-attention, you normally compute something like:
\[
  \operatorname{Attention}(Q,K,V)
  =\operatorname{softmax}\!\left(\frac{QK^{\mathsf T}}{\sqrt{d}}\right)V.
\]
The problem is that plain attention by itself does not know whether a token
is at position 3, 30, or 3000.

RoPE solves that by rotating the query and key vectors by an angle that
depends on token position. Conceptually:
\begin{itemize}
    \setlength{\itemsep}{0pt}
  \item token at position 0 $\rightarrow$ rotate Q/K by angle 0;
  \item token at position 1 $\rightarrow$ rotate Q/K a little;
  \item token at position 2 $\rightarrow$ rotate Q/K more;
  \item token at position 3 $\rightarrow$ rotate Q/K more again;
  \item \ldots
\end{itemize}
The clever part is that when the rotated query and key interact through the
dot product, their relationship naturally contains information about the
distance between positions. So the model can learn things like ``this word
is 1 token behind me,'' ``this variable was defined 40 tokens ago,'' or
``this sentence started 200 tokens earlier'' --- rather than merely learning
an absolute position number.

A simplified 2D rotation looks like this:
\[
  \begin{bmatrix} x_1' \\ x_2'
  \end{bmatrix}
  =
  \begin{bmatrix}
    \cos\theta & -\sin\theta \\
    \sin\theta & \cos\theta
  \end{bmatrix}
  \begin{bmatrix} x_1 \\ x_2
  \end{bmatrix},
\]
where \(\theta\) depends on token position and dimension/frequency. In
practice, RoPE applies these rotations across pairs of dimensions in \(Q\) and
\(K\).

\subsubsection{Simple example}

Suppose one attention head has only 4 dimensions:
\[
  q=[q_1,q_2,q_3,q_4].
\]
RoPE treats these as two 2D pairs \((q_1,q_2)\) and \((q_3,q_4)\), and rotates
each pair according to the token's position.

RoPE is not primarily useful because the model knows ``I am at position 17.''
The really useful property appears when \(Q\) interacts with \(K\).

Suppose we simplify further and use \(q=k=[1,0]\) before RoPE. \\
Imagine the query is at position 3 and the key is at position 1, with
\(30^{\circ}\) per
position. Then:
\[
  Q_3=[\cos 90^{\circ},\sin 90^{\circ}]=[0,1],
\]
\[
  K_1=[\cos 30^{\circ},\sin 30^{\circ}]=[0.866,0.5].
\]
Attention uses their dot product:
\[
  Q_3\cdot K_1=(0)(0.866)+(1)(0.5)=0.5.
\]
The positional difference is also: \\
\((3-1)\times 30^{\circ}=60^{\circ}\), and
\(\cos 60^{\circ}=0.5\).

There is a trigonometric identity underneath: if one vector is rotated by
angle \(a\) and another by angle \(b\), their relative angle becomes
\(a-b\). For RoPE,
\[
  a=p_q\,\theta,\qquad b=p_k\,\theta,
\]
so the important term becomes \((p_q-p_k)\theta\). \\
Therefore the attention
interaction contains \(p_q-p_k\) --- the distance between the two tokens.

\paragraph{Why multiple frequencies?}
Our 4D example contained two pairs, \((q_1,q_2)\) and \((q_3,q_4)\), rotated
at different speeds. Modern RoPE does this across many dimensions:
\begin{itemize}
    \setlength{\itemsep}{0pt}
  \item dimension pair 0 $\rightarrow$ relatively high frequency;
  \item dimension pair 1 $\rightarrow$ slightly lower frequency;
  \item dimension pair 2 $\rightarrow$ lower;
  \item dimension pair 3 $\rightarrow$ lower;
  \item \ldots
\end{itemize}

\texttt{rope\_cache()} precomputes sine and cosine tables of shape
\([T,D/2]\). \\
\texttt{apply\_rope()} splits each query or key head into two
halves and applies a position-dependent rotation:

\[
  (x_1,x_2)\longmapsto
  (x_1\cos\theta-x_2\sin\theta,
  x_1\sin\theta+x_2\cos\theta).
\]

Values are not rotated. Position affects which keys a query aligns with, while
the values continue to carry content.

\subsection{Grouped-query causal attention (GQA)}

Previously we studied normal multi-head attention where:
\begin{lstlisting}[style=shadedpython,language={},basicstyle=\scriptsize\ttfamily]
8 query heads
8 key heads
8 value heads
\end{lstlisting}
In grouped-query attention (GQA), key/value heads are grouped.

Conceptually:
\begin{lstlisting}[style=shadedpython,language={},basicstyle=\scriptsize\ttfamily]
Q heads

Q1 Q2 Q3 Q4 Q5 Q6 Q7 Q8

      share KV groups

K1/V1 serves Q1-Q4
K2/V2 serves Q5-Q8
\end{lstlisting}
because
\begin{lstlisting}[style=shadedpython]
rep = n_heads // n_kv_heads
\end{lstlisting}
Here \(8/2=4\), so each KV head serves 4 query heads. This reduces KV memory
and computation.

Queries are projected to \(H\) heads, while keys and values are projected to
only \(H_{kv}\) heads:

\[
  Q:[B,H,T,D],\qquad
  K,V:[B,H_{kv},T,D].
\]

Each KV head serves \(H/H_{kv}\) query heads. The code currently realizes this
by \texttt{repeat\_interleave} before calling PyTorch scaled-dot-product
attention. This saves KV projection parameters, but materializing repeated K/V
tensors gives up much of GQA's intended activation and cache-memory benefit.
A production kernel should consume grouped K/V directly.

\texttt{is\_causal=True} delegates causal masking and fused attention-kernel
selection to PyTorch. The result is concatenated across heads and projected
back to \(C\).

\subsection{SwiGLU and the residual block}

SwiGLU is basically a smarter feed-forward network (FFN) used inside
many modern Transformers.

A traditional FFN looks like:
\begin{lstlisting}[style=shadedpython,language={},basicstyle=\scriptsize\ttfamily]
x
|
v
Linear up-projection
|
v
GELU / ReLU
|
v
Linear down-projection
|
v
output
\end{lstlisting}
Mathematically, roughly:
\[
  \operatorname{FFN}(x)=W_{down}\bigl(\operatorname{GELU}(W_{up}x)\bigr).
\]

\subsubsection{SwiGLU introduces two branches}

Instead of making one projection from \(x\), SwiGLU makes two:
\begin{lstlisting}[style=shadedpython,language={},basicstyle=\scriptsize\ttfamily]
                   x
                 /   \
                /     \
               v       v
          W_gate x    W_up x
               |       |
             SiLU      |
               \       /
                \     /
                  *
                  |
                  v
               W_down
                  |
                  v
                output
\end{lstlisting}
So \(x\) goes through two independent learned matrices:
\[
  g=W_{gate}x,\qquad u=W_{up}x.
\]
The first becomes the gate; the second contains the features/values.

\subsubsection{The gate goes through SiLU}

We apply \(\operatorname{SiLU}(g)\), where
\[
  \operatorname{SiLU}(x)=x\,\sigma(x)
\]
and \(\sigma(x)\) is the sigmoid. Conceptually SiLU behaves somewhat like:
\begin{center}
  \renewcommand{\arraystretch}{1.15}
  \small
  \begin{tabular}{@{}r r@{}}
    \textbf{Input} & \textbf{SiLU output} \\
    \hline
    \(-5\) & near \(0\) \\
    \(-1\) & about \(-0.27\) \\
    \(0\) & \(0\) \\
    \(1\) & about \(0.73\) \\
    \(5\) & about \(4.97\)
  \end{tabular}
\end{center}
So it smoothly controls how strongly something passes through. That is why we
call it a gate.

\subsubsection{Then multiply the two branches}

This is the crucial bit. Suppose
\[
  W_{gate}x=[2,-2,0.5,1]
\]
and
\[
  W_{up}x=[10,20,30,40].
\]
After SiLU, imagine approximately
\[
  \operatorname{SiLU}(W_{gate}x)=[1.76,-0.24,0.31,0.73].
\]
Now multiply element-by-element:
\[
  [1.76,-0.24,0.31,0.73]\odot[10,20,30,40]
\]
giving approximately \([17.6,-4.8,9.3,29.2]\). The gate has therefore altered
how much of every feature gets through. You can mentally read that as:
\begin{itemize}
    \setlength{\itemsep}{0pt}
  \item Feature 1 $\rightarrow$ strongly useful $\rightarrow$ let it through;
  \item Feature 2 $\rightarrow$ suppress/change it;
  \item Feature 3 $\rightarrow$ somewhat useful;
  \item Feature 4 $\rightarrow$ quite useful.
\end{itemize}

\subsubsection{Then \(W_{down}\) brings it back}

The intermediate SwiGLU vector is usually much larger than the model's hidden
dimension, so \(W_{down}\) projects it back down. The complete equation is:
\[
  \operatorname{SwiGLU}(x)
  =W_{down}\bigl(\operatorname{SiLU}(W_{gate}x)\odot W_{up}x\bigr),
\]
where \(\odot\) means element-wise multiplication.

\subsubsection{Why is it called SwiGLU?}

The name can be broken apart. \emph{Swish} is essentially the activation now
commonly called SiLU:
\[
  \operatorname{SiLU}(x)=x\,\sigma(x).
\]
And \emph{GLU} means Gated Linear Unit.

\subsection{Transformer block}

So inside your eventual Transformer block:
\begin{lstlisting}[style=shadedpython,language={},basicstyle=\scriptsize\ttfamily]
                 INPUT
                   |
              +----+----+
              |Attention|  <- RoPE operates in Q/K here
              +----+----+
                   |
               residual
                   |
              +----+----+
              | SwiGLU  |  <- gated feature processing here
              |   FFN   |
              +----+----+
                   |
               residual
                   |
                 OUTPUT
\end{lstlisting}

The block is:
\begin{lstlisting}[style=shadedpython]
x = x + Attention(RMSNorm(x))
x = x + SwiGLU(RMSNorm(x))
\end{lstlisting}

\subsection{Embedding, head, loss, and generation}

Token embeddings map \([B,T]\) IDs to \([B,T,C]\). After all blocks and a final
RMSNorm, the language-model head maps to \([B,T,V]\). Its weight is the exact
same parameter as the token-embedding table, reducing parameters and coupling
input and output token geometry.

Training flattens logits and targets for cross-entropy. Inference without
targets computes only the last position's logits. \texttt{generate()} repeatedly
crops context to the maximum sequence length, samples from softmax, and appends
one token. It has no KV cache, top-k/top-p control, EOS stopping, or batched
finished-sequence handling; it is a correctness path, not an inference engine.

\subsection{Training \texttt{train.py}}

\begin{lstlisting}[style=shadedpython,language={},basicstyle=\scriptsize\ttfamily]
train.bin / val.bin
       |
       v
train.py
       |
       v
Data
       |
       v
batches
       |
       v
GPT
       |
       v
loss
       |
       v
backprop
       |
       v
AdamW
       |
       v
checkpoint
\end{lstlisting}

\subsection{Memory-mapped training data}

\begin{lstlisting}[style=shadedpython]
np.memmap(...)
\end{lstlisting}
is important. Instead of loading 1.3B tokens into RAM, the operating system
maps the binary file. So the trainer can access slices like
\begin{lstlisting}[style=shadedpython]
arr[s:s+1024]
\end{lstlisting}
without loading the full dataset.

\subsection{Prepared data and deterministic batches}

The \texttt{Data} class memory-maps flat token arrays. For a training micro-step
it seeds NumPy with the tuple

\[
  (\text{seed},\text{optimizer step},\text{micro-step},\text{is validation})
\]

and samples start offsets. Inputs and one-token-shifted targets are converted to
\texttt{int64} PyTorch tensors and moved to CUDA, MPS, or CPU. Thus
the batch for
a particular step is independent of earlier process history.

Validation intentionally reuses the same fixed batches, making measurements
comparable across checkpoints. A later distributed version must incorporate
rank into training seeds without allowing ranks to sample identical sequences.

\subsection{Optimizer and schedule}

Parameters with at least two dimensions receive AdamW weight decay; 1-D norm
weights do not. The learning rate warms linearly and then follows cosine decay
to 10\% of its peak by default.

For each optimizer step, the loop performs:

\begin{enumerate}
    \setlength{\itemsep}{0pt}
  \item set the learning rate from the global step;
  \item clear gradients;
  \item run one or more deterministic micro-batches;
  \item divide every micro-loss by the accumulation count;
  \item backpropagate and accumulate gradients;
  \item clip the global gradient norm; and
  \item call AdamW once.
\end{enumerate}

CUDA uses BF16 autocast. MPS and CPU run FP32. There is no FP16 gradient scaler,
which is consistent with the current precision choices.

\subsection{Checkpoint semantics}

It saves:
\begin{lstlisting}[style=shadedpython]
{
    "model": ...,
    "opt": ...,
    "step": ...,
    "model_config": ...,
    "train_config": ...
}
\end{lstlisting}
So a checkpoint includes model weights, optimizer state, current step, model
architecture, and training configuration.

With
\begin{lstlisting}[style=shadedpython]
python train.py \
    --config ... \
    --resume
\end{lstlisting}
training can continue from \texttt{ckpt\_last.pt}.

The checkpoint stores model weights, AdamW state, completed step, and both
configuration dictionaries. Resumption loads them and begins at
\texttt{saved step + 1}. Learning rate is a pure function of the step, and
training batches are a pure function of step and seed. Because the model has no
dropout, this is sufficient for the present deterministic smoke drill.

It is not yet a robust distributed checkpoint. The write is not atomic; random
states are not stored; there is only one latest file; dataloader, scaler, and
distributed-shard state do not exist; and corrupt-checkpoint fallback is absent.

\subsection{Logging and MFU}

Tokens per optimizer step are

\[
  B\times T\times\text{gradient accumulation}.
\]

The trainer approximates training work as six times parameter count per token
and divides estimated FLOP/s by configured device peak FLOP/s. This is a useful
rough MFU signal, not a benchmark-quality measurement: it omits quadratic
attention work and communication, and the interval timing logic needs care for
resumed or partial logging windows.

\subsection{What the configurations actually represent}

Two different quantities get confused here: \emph{parameters} and
\emph{tokens}.

\begin{itemize}
  \setlength{\itemsep}{0pt}
  \item \textbf{Parameters} count the trainable numbers in the model
    (weights). We write M for million (\(10^6\)), B for billion
    (\(10^9\)). So ``21.30M parameters'' means about \(21.3\times 10^6\)
    weights.
  \item \textbf{Tokens} count how much training text the optimizer sees.
    T means trillion (\(10^{12}\)). So ``train on 2T tokens'' means feed
    the model about \(2\times 10^{12}\) BPE tokens --- it does \emph{not}
    mean the model has 2 trillion weights.
\end{itemize}

A small model can still train on a large token budget (many passes over
data, or a huge corpus). Conversely, a large-parameter model can be
under-trained if it sees too few tokens. The corpus roadmap's ``2T'' target
is a \emph{data} scale (tokens), not a parameter scale.

\begin{center}
  \renewcommand{\arraystretch}{1.2}
  \begin{tabular}{@{}l r r r r@{}}
    \textbf{Config} & \(C\) & \(L\) & \(H/H_{kv}\) & \textbf{Parameters} \\
    \hline
    smoke & 128 & 4 & 4/2 & 1.21M with actual 1{,}745 vocab \\
    mini & 256 & 6 & 8/2 & 21.30M \\
    300m & 1{,}024 & 24 & 16/4 & 337.69M \\
    1b & 2{,}048 & 26 & 16/4 & 1.307B
  \end{tabular}
\end{center}

The names are scale labels, not exact counts. The repository does not yet have
30--50M or 100--150M configurations.

\subsubsection{How 21.30M is calculated (\texttt{mini.json})}

The code reports
\[
  \texttt{num\_params()}
  =\sum_p |p|,
\]
i.e.\ the sum of every tensor's element count (with the LM head tied to the
token embedding, so that table is counted once). For
\texttt{mini.json} (\(V{=}65{,}536\), \(C{=}256\), \(L{=}6\),
\(H{=}8\), \(H_{kv}{=}2\)):

\begin{center}
  \renewcommand{\arraystretch}{1.15}
  \small
  \begin{tabular}{@{}p{0.42\textwidth}r@{}}
    \textbf{Piece} & \textbf{Count} \\
    \hline
    Token embedding \(V\times C\)
      (tied LM head) &
    \(65{,}536\times 256=16{,}777{,}216\) \\
    Attention per layer
      (\(W_q,W_k,W_v,W_o\); GQA shrinks \(K/V\)) &
    \(163{,}840\) \\
    SwiGLU per layer
      (hidden \(=768\) after \(8C/3\) rounded up) &
    \(589{,}824\) \\
    RMSNorm weights per layer (\(2\times C\)) &
    \(512\) \\
    \(\times\,L{=}6\) Transformer blocks &
    \(6\times 754{,}176=4{,}525{,}056\) \\
    Final RMSNorm & \(256\) \\
    \hline
    \textbf{Total} &
    \(21{,}302{,}528\approx 21.30\text{M}\)
  \end{tabular}
\end{center}

Notice the embedding alone is \(\sim\)79\% of this tiny model. That is why
scaling width/depth (not only vocab) is what grows ``real'' Transformer
capacity.

\subsection{Required increments before larger runs}

\begin{enumerate}
    \setlength{\itemsep}{0pt}
  \item Add unit tests for shapes, causal behavior, tied weights,
    RoPE, and resume equivalence.
  \item Make validation document-disjoint and domain-balanced.
  \item Add atomic checkpoint writes, retained milestones, and recovery tests.
  \item Add \texttt{torch.compile}, fused AdamW where safe, and profiler traces.
  \item Implement DDP for throughput and FSDP or an adopted framework for 1B.
  \item Add gradient checkpointing and a direct GQA attention path.
  \item Log structured JSON with tokens seen, validation loss,
    gradient norm, memory, and config hashes.
  \item Reconcile every run with a snapshot ID and experiment-registry entry.
\end{enumerate}
